\documentclass[10pt]{article}
\usepackage[T1]{fontenc}
\usepackage[utf8]{inputenc}
\usepackage[french]{babel}
\usepackage{amsmath,amssymb,mathrsfs,mathpazo}
\usepackage{geometry}
\geometry{a4paper,margin=2cm}
\usepackage{graphicx,array,tabularx,tabularray,colortbl,multirow,multicol,booktabs,enumitem,pifont,ulem,fancybox,diagbox,cancel}
\usepackage{pstricks,pstricks-add,pst-plot,pst-tree,pst-node,pst-text,pst-func,pst-math,pst-bspline,pst-3dplot,pst-eucl,pst-solides3d}
\let\pscalculate\undefined
\usepackage{pst-bezier}
\usepackage{tikz,pgfplots,tkz-tab}
\usetikzlibrary{arrows,arrows.meta,patterns,calc,positioning,trees,decorations.markings,shapes.geometric,shapes.misc,shapes.arrows,intersections,decorations.pathmorphing,decorations.pathreplacing,backgrounds,shadows,babel}
\pgfplotsset{compat=1.18}
\usepgfplotslibrary{fillbetween}
\usepackage[np]{numprint}
\usepackage{listings,makecell,pas-tableur,twoopt}
\let\barre\undefined
\usepackage[pstricks]{tablvar}
\definecolor{darkgreen}{RGB}{0,100,0}
\providecommand{\red}{\color{red}}
\providecommand{\blue}{\color{blue}}
\providecommand{\green}{\color{green}}
\providecommand{\black}{\color{black}}
\providecommand{\white}{\color{white}}
\providecommand{\gray}{\color{gray}}
\providecommand{\cyan}{\color{cyan}}
\providecommand{\magenta}{\color{magenta}}
\providecommand{\yellow}{\color{yellow}}
\providecommand{\violet}{\color{violet}}
\providecommand{\orange}{\color{orange}}
\providecommand{\euro}{\texteuro}
\providecommand{\EUR}{\texteuro}
\providecommand{\ieme}{\textsuperscript{e}}
\providecommand{\degres}{\ensuremath{^\circ}}
\providecommand{\degree}{\ensuremath{^\circ}}
\providecommand{\vv}[1]{\overrightarrow{#1}}
\providecommand{\Coord}[1]{#1}
\providecommand{\decosix}{}
\providecommand{\starredbullet}{\ensuremath{\bullet}}
\pagestyle{empty}
\setlength{\parindent}{0pt}
\usepackage{adjustbox}\DeclareUnicodeCharacter{2236}{:}\DeclareUnicodeCharacter{2190}{\ensuremath{\leftarrow}}
\providecommand{\R}{\mathbb{R}}
\providecommand{\N}{\mathbb{N}}
\providecommand{\Z}{\mathbb{Z}}
\providecommand{\Q}{\mathbb{Q}}
\providecommand{\D}{\mathbb{D}}
\providecommand{\C}{\mathbb{C}}
\providecommand{\e}{{\mathrm{e}}}
\providecommand{\vect}[1]{\overrightarrow{#1}}
\providecommand{\vectt}[1]{\overrightarrow{\mathrm{#1}}}
\providecommand{\vvt}[1]{\overrightarrow{\mathrm{#1}}}
\providecommand{\np}[2][]{{#2~#1}}
\providecommand{\vv}[1]{{\overrightarrow{#1}}}
\renewcommand{\d}{{\mathrm{d}}}
\providecommand{\degres}{{^\circ}}
\providecommand{\degre}{{^\circ}}
\providecommand{\blue}{\color{blue}}
\providecommand{\red}{\color{red}}
\providecommand{\enskip}{\quad}
\providecommand{\llbracket}{\lbrack\!\lbrack}
\providecommand{\rrbracket}{\rbrack\!\rbrack}
\providecommand{\euro}{{\text{€}}}
\providecommand{\no}{{\text{n°}}}
\providecommand{\interval}[2]{{[#1;#2]}}
\providecommand{\intervFO}[2]{{[#1;#2[}}
\providecommand{\dotfill}{\ldots}

\providecommand{\ds}{\displaystyle}
\providecommand{\textding}[1]{#1}
\providecommand{\Oij}{\(\mathrm{O};\vec{\imath},\vec{\jmath}\)}
\providecommand{\Oijk}{\(\mathrm{O};\vec{\imath},\vec{\jmath},\vec{k}\)}
\providecommand{\Ouv}{\(\mathrm{O};\vec{u},\vec{v}\)}
\providecommand{\intervFF}[2]{\ensuremath{[#1;#2]}}
\renewcommand{\euro}{\ifmmode\text{\texteuro}\else\texteuro\fi}
\DeclareUnicodeCharacter{00B0}{\ifmmode^{\circ}\else\textdegree\fi}
\DeclareUnicodeCharacter{20AC}{\ifmmode\text{\texteuro}\else\texteuro\fi}
% Énoncés : Éducation nationale. Transcription et corrigés : APMEP, Denis Vergès et contributeurs. Découpage et classement : CoopMaths, Sébastien Lozano, Éric Elter et contributeurs. Adaptation de la présentation : Maths Explorer.
% CC BY-SA : https://coopmaths.fr/www/about/
% https://coopmaths.fr/alea/static/bac/2024/tex/bac_2024_01_sujet0_8_cor.tex
\begin{document}
\noindent\textit{APMEP / CoopMaths. CC BY-SA. Présentation adaptée par Maths Explorer.}\par\medskip
On considère la suite $\left(u_{n}\right)$ définie par $u_{0} = 0$ et $u_{n + 1} =  3 u_{n} + 1$ pour tout entier naturel $n$.

\begin{enumerate}
	\item 
\setlength{\columnsep}{1cm} 
\setlength{\columnseprule}{0.5pt} 
\begin{multicols}{2}	
	On considère la fonction \texttt{calcul} écrite dans le langage Python qui renvoie la valeur de $u_{n}$.
	
\begin{center}
\renewcommand\arraystretch{1}
\texttt{
\begin{tabular}{l}
def calcul(n):\\
\quad u = 0\\
\quad for i in range(n):\\
\quad\quad u = 3 * u + 1\\
\quad return u\\
\end{tabular}}
\end{center}

\columnbreak

On considère par ailleurs la fonction \texttt{liste} écrite dans le langage Python :

\begin{center}
\renewcommand\arraystretch{1}
\texttt{
		\begin{tabular}{l}
		def liste(n):\\
		\quad l = []\\
		\quad for i in range(n):\\
		\quad\quad l.append(calcul(i))\\
		\quad return l\\
		\end{tabular}}
\end{center}

\end{multicols}

	\textbf{Affirmation 1 :} \og{}l'appel \texttt{liste(6)} renvoie la liste \texttt{[0, 1, 4, 13, 42, 121]}.\fg{}
	
$u_0=0$, $u_1=3\times u_0+1=1$, $u_2=3\times u_1+1=4$, $u_3=3\times u_2+1=13$ et\\
 $u_4=3\times u_3+1=40\neq 42$.

\hfill \textbf{Affirmation 1 fausse}
	
	\item \textbf{Affirmation 2 :} \og{} pour tout entier naturel $n$, $u_{n} = \dfrac{1}{2} \times 3^{n}-\dfrac{1}{2}$.\fg{}

Pour $n=0$, on a $u_0=0$ et $\dfrac{1}{2} \times 3^{0}-\dfrac{1}{2} = \dfrac{1}{2}-\dfrac{1}{2}=0$.

Pour $n=1$, on a $u_1=1$ et $\dfrac{1}{2} \times 3^{1}-\dfrac{1}{2} = \dfrac{3}{2}-\dfrac{1}{2}=1$.

Pour $n=2$, on a $u_2=4$ et $\dfrac{1}{2} \times 3^{2}-\dfrac{1}{2} = \dfrac{9}{2}-\dfrac{1}{2}=4$.

On conjecture que la propriété \og{} $u_{n} = \dfrac{1}{2} \times 3^{n}-\dfrac{1}{2}$\fg{} est vraie pour tout $n$, et on va démontrer cette conjecture en utilisant un raisonnement par récurrence.

\begin{list}{\textbullet}{}
\item \textbf{Initialisation}

Pour $n=0$, on a $u_0=0$ et $\dfrac{1}{2} \times 3^{0}-\dfrac{1}{2} = \dfrac{1}{2}-\dfrac{1}{2}=0$; la propriété est vraie.

\item \textbf{Hérédité}

On suppose la propriété vraie au rang $n$, c'est-à-dire $u_{n} = \dfrac{1}{2} \times 3^{n}-\dfrac{1}{2}$.

$u_{n+1}=3u_n+1 = 3\left (\dfrac{1}{2} \times 3^{n}-\dfrac{1}{2}\right )+1
= \dfrac{1}{2}\times 3^{n+1} -\dfrac{3}{2}+1
= \dfrac{1}{2}\times 3^{n+1} -\dfrac{1}{2}$

La propriété est donc vraie au rang $n+1$.

\item \textbf{Conclusion}

La propriété est vraie au rang $0$, et elle est héréditaire pour tout $n\geqslant 0$, donc, d'après le principe de récurrence, elle est vraie pour tout $n$.
\end{list}

\hfill \textbf{Affirmation 2 vraie}

\item \textbf{Affirmation 3 :} \og{} pour tout entier naturel $n$, \quad $u_{n + 1}-u_{n}$ est une puissance de 3.\fg{}

$\begin{gathered}u_{n+1}-u_n = \left (\dfrac{1}{2}\times 3^{n+1} -\dfrac{1}{2}\right ) - \left ( \dfrac{1}{2}\times 3^{n} -\dfrac{1}{2}\right )
= \dfrac{1}{2}\times 3^{n+1} -\dfrac{1}{2} - \dfrac{1}{2}\times 3^{n} +\dfrac{1}{2}\\
\phantom{u_{n+1}-u_n}
= \dfrac{1}{2}\times 3^{n+1} - \dfrac{1}{2}\times 3^{n}
= \dfrac{1}{2}\times 3^n\left (3-1\right )
= \dfrac{1}{2}\times 3^n \times 2
= 3^n\end{gathered}$
\hfill \textbf{Affirmation 3 vraie}
\end{enumerate}

\end{document}
